\paragrafos{Ισοδύναμα κλάσματα: το ίδιο κομμάτι με διαφορετικό κόψιμο}{Α.2.2, «Ισοδύναμα κλάσματα»}

\begin{anagke}
Παρασκευή βράδυ στο σπίτι του Βασίλη, που γιορτάζει τη νίκη του στις εκλογές της τάξης με πίτσες. Οι πίτσες είναι τέσσερις, ίδιου μεγέθους, αλλά κάθε παιδί κόβει τη δική του όπως θέλει. Ο Νίκος τη χωρίζει σε 4 κομμάτια, η Αγγελική σε 6, ο Βασίλης σε 8 και η Ζωή σε 10. Και οι τέσσερις τρώνε κάποια κομμάτια και ανακοινώνουν: ο Νίκος 2, η Αγγελική 3, ο Βασίλης 4 και η Ζωή 5.

«Εγώ έφαγα δύο, εσύ τρία, εγώ τέσσερα, εσύ πέντε. Ο Βασίλης έφαγε τα πιο πολλά!» λέει θριαμβευτικά ο Στέλιος, που μετράει μόνο κομμάτια. Η Πολυξένη κοιτάζει τις τέσσερις πίτσες πάνω στο τραπέζι και λέει ήσυχα: «Έφαγαν όλοι το ίδιο. Όλοι έφαγαν τη μισή πίτσα».

«Μα ο ένας έφαγε δύο και ο άλλος πέντε!» διαμαρτύρεται ο Στέλιος. Η Μυρτώ σηκώνει το κομμάτι της Ζωής και το βάζει δίπλα στης Αγγελικής. «Δεν είναι ίδια κομμάτια. Τα κομμάτια της Ζωής είναι μικρότερα». Πώς γράφουμε με αριθμούς το ότι τα $\frac{2}{4}$, τα $\frac{3}{6}$, τα $\frac{4}{8}$ και τα $\frac{5}{10}$ είναι το ίδιο κομμάτι;
\end{anagke}

\begin{lexi}{ισοδύναμα, απλοποίηση, ανάγωγο}{ισοδύναμα, απλοποίηση, ανάγωγο}
Δύο κλάσματα είναι \lexiword{ισοδύναμα} όταν έχουν την ίδια «δύναμη», δηλαδή την ίδια αξία. Το $\frac{2}{4}$ και το $\frac{3}{6}$ δεν είναι ίδια σύμβολα, αλλά είναι ίσα, γιατί δείχνουν το ίδιο κομμάτι.

\lexiword{Απλοποιώ} ένα κλάσμα σημαίνει το κάνω πιο απλό, με μικρότερους αριθμούς, χωρίς να αλλάξει η αξία του. Το \lexiword{ανάγωγο} κλάσμα είναι αυτό που δεν μπορεί να γίνει πιο απλό: δεν μπορεί να «αναχθεί», δηλαδή να φερθεί πίσω σε μικρότερους όρους.

Και δύο λέξεις για τα ονόματα: \lexiword{ομώνυμα} είναι τα κλάσματα που έχουν το ίδιο όνομα (παρονομαστή), και \lexiword{ετερώνυμα} εκείνα που έχουν διαφορετικό.
\end{lexi}

\begin{metafrasi}
Η μισή πίτσα, κομμένη σε 4 ή σε 6 κομμάτια. & $\dfrac{2}{4}=\dfrac{3}{6}$\\
Τα χιαστί γινόμενα είναι ίσα. & $2\cdot 6=4\cdot 3$\\
Πολλαπλασιάζω και τους δύο όρους επί 3. & $\dfrac{2}{5}=\dfrac{2\cdot 3}{5\cdot 3}=\dfrac{6}{15}$\\
Διαιρώ και τους δύο όρους με το 14. & $\dfrac{42}{56}=\dfrac{42:14}{56:14}=\dfrac{3}{4}$\\
Δεν απλοποιείται άλλο (ΜΚΔ $=1$). & $\dfrac{7}{12}$ είναι ανάγωγο\\
\end{metafrasi}

\tmima{Η θεωρία}

\begin{orismos}{ισοδύναμα κλάσματα}
Δύο κλάσματα $\dfrac{\alpha}{\beta}$ και $\dfrac{\gamma}{\delta}$ λέγονται \textbf{ισοδύναμα} ή \textbf{ίσα} όταν εκφράζουν το ίδιο μέρος ενός μεγέθους. Γράφουμε $\dfrac{\alpha}{\beta}=\dfrac{\gamma}{\delta}$.
\end{orismos}

\begin{idiotita}{χιαστί γινόμενα}
Αν $\dfrac{\alpha}{\beta}=\dfrac{\gamma}{\delta}$, τότε τα «χιαστί γινόμενα» είναι ίσα: $\alpha\cdot\delta=\beta\cdot\gamma$. Και αντιστρόφως.
\end{idiotita}

Έτσι $\frac{3}{4}=\frac{9}{12}$, γιατί $3\cdot 12=36=4\cdot 9$. Και τα $\frac{3}{5}$ και $\frac{4}{7}$ δεν είναι ισοδύναμα, γιατί $3\cdot 7=21$ ενώ $5\cdot 4=20$.

\begin{kanonas}{ισοδύναμα κλάσματα}
\begin{itemize}[leftmargin=1.2em,itemsep=0.3ex,topsep=2pt]
\item Αν πολλαπλασιάσω \textbf{και τους δύο} όρους ενός κλάσματος με τον ίδιο φυσικό αριθμό (όχι 0), βγαίνει ισοδύναμο κλάσμα: $\dfrac{2}{3}=\dfrac{2\cdot 4}{3\cdot 4}=\dfrac{8}{12}$.
\item Αν διαιρέσω \textbf{και τους δύο} όρους με τον ίδιο φυσικό αριθμό (όχι 0, που διαιρεί και τους δύο), βγαίνει ισοδύναμο κλάσμα: $\dfrac{18}{30}=\dfrac{18:6}{30:6}=\dfrac{3}{5}$.
\end{itemize}
\end{kanonas}

Η δεύτερη διαδικασία λέγεται \textbf{απλοποίηση}. Αν διαιρέσουμε και τους δύο όρους με τον \textbf{ΜΚΔ} τους, το κλάσμα γίνεται ανάγωγο σε ένα βήμα. Ένα κλάσμα λέγεται \textbf{ανάγωγο} όταν δεν έχει κοινό διαιρέτη ο αριθμητής και ο παρονομαστής εκτός από το 1, δηλαδή όταν είναι πρώτοι μεταξύ τους. Το $\frac{7}{12}$ είναι ανάγωγο, γιατί $\MKD(7,\,12)=1$.

\tmimasmall{Ομώνυμα και ετερώνυμα}

Τα κλάσματα που έχουν τον ίδιο παρονομαστή λέγονται \textbf{ομώνυμα}, όπως τα $\frac{2}{7}$ και $\frac{5}{7}$. Αυτά που έχουν διαφορετικούς παρονομαστές λέγονται \textbf{ετερώνυμα}, όπως τα $\frac{2}{7}$ και $\frac{5}{3}$.

\begin{kanonas}{πώς κάνω ετερώνυμα κλάσματα ομώνυμα}
\begin{enumerate}[leftmargin=1.6em,itemsep=0.2ex,topsep=2pt]
\item Απλοποιώ όσα κλάσματα απλοποιούνται.
\item Βρίσκω το \textbf{ΕΚΠ} των παρονομαστών: θα είναι ο κοινός παρονομαστής.
\item Διαιρώ το ΕΚΠ με τον παρονομαστή κάθε κλάσματος.
\item Πολλαπλασιάζω και τους δύο όρους κάθε κλάσματος με τον αριθμό που βρήκα για εκείνο το κλάσμα.
\end{enumerate}
\end{kanonas}

\begin{prosoxi}
\begin{itemize}[leftmargin=1.2em,itemsep=0.3ex,topsep=2pt]
\item Πολλαπλασιάζουμε ή διαιρούμε \textbf{και τους δύο} όρους με τον \textbf{ίδιο} αριθμό. Αν προσθέσεις τον ίδιο αριθμό και στους δύο όρους, αλλάζει η αξία: το $\frac{1}{2}$ δεν είναι ίσο με το $\frac{1+3}{2+3}=\frac{4}{5}$.
\item Στην απλοποίηση ο αριθμός που διαιρεί πρέπει να διαιρεί \textbf{και τους δύο} όρους ακριβώς.
\item Αν πολλαπλασιάσεις με το 0, χάνεις το κλάσμα: δεν επιτρέπεται.
\end{itemize}
\end{prosoxi}

\tmima{Λυμένα παραδείγματα}

\paradeigma{Οι πίτσες του Βασίλη}
\tixero Ο Νίκος έφαγε $\frac{2}{4}$ της πίτσας του, η Αγγελική $\frac{3}{6}$ και ο Βασίλης $\frac{4}{8}$.\par
\tithelo Να εξετάσω αν έφαγαν το ίδιο.\par
\possindeo Υπολογίζω τα χιαστί γινόμενα. Για το $\frac{2}{4}$ και το $\frac{3}{6}$: $2\cdot 6=12$ και $4\cdot 3=12$. Ίσα, άρα ισοδύναμα. Για το $\frac{3}{6}$ και το $\frac{4}{8}$: $3\cdot 8=24$ και $6\cdot 4=24$. Ίσα, άρα ισοδύναμα.\par
\apantisi Έφαγαν το ίδιο: όλοι τη μισή πίτσα.\par
\elegxos Και τα τρία απλοποιούνται στο $\frac{1}{2}$: $\frac{2}{4}=\frac{2:2}{4:2}=\frac{1}{2}$, $\frac{3}{6}=\frac{3:3}{6:3}=\frac{1}{2}$, $\frac{4}{8}=\frac{4:4}{8:4}=\frac{1}{2}$.

\paradeigma{Απλοποίηση του $\frac{84}{126}$}
\tixero Το κλάσμα $\frac{84}{126}$.\par
\tithelo Να το κάνω ανάγωγο.\par
\possindeo Βρίσκω το ΜΚΔ των όρων. $84=2^2\cdot 3\cdot 7$ και $126=2\cdot 3^2\cdot 7$, άρα $\MKD(84,\,126)=2\cdot 3\cdot 7=42$. Διαιρώ και τους δύο όρους με το 42: $\frac{84:42}{126:42}=\frac{2}{3}$.\par
\apantisi $\dfrac{84}{126}=\dfrac{2}{3}$.\par
\elegxos Το $\frac{2}{3}$ είναι ανάγωγο, γιατί το 2 και το 3 είναι πρώτοι μεταξύ τους. Και τα χιαστί γινόμενα: $84\cdot 3=252=126\cdot 2$. Σωστά.

\paradeigma{Τρία κλάσματα, ένας παρονομαστής}
\tixero Τα κλάσματα $\frac{3}{4}$, $\frac{5}{6}$ και $\frac{7}{10}$.\par
\tithelo Να τα κάνω ομώνυμα.\par
\possindeo Κανένα δεν απλοποιείται, γιατί είναι όλα ανάγωγα. Βρίσκω το $\EKP(4,\,6,\,10)$: $4=2^2$, $6=2\cdot 3$, $10=2\cdot 5$, άρα $2^2\cdot 3\cdot 5=60$. Διαιρώ: $60:4=15$, $60:6=10$, $60:10=6$. Πολλαπλασιάζω τους όρους: $\frac{3}{4}=\frac{3\cdot 15}{4\cdot 15}=\frac{45}{60}$, $\frac{5}{6}=\frac{5\cdot 10}{6\cdot 10}=\frac{50}{60}$, $\frac{7}{10}=\frac{7\cdot 6}{10\cdot 6}=\frac{42}{60}$.\par
\apantisi $\dfrac{45}{60}$, $\dfrac{50}{60}$ και $\dfrac{42}{60}$.\par
\elegxos Απλοποιώντας το $\frac{45}{60}$ με το 15 ξαναβρίσκω το $\frac{3}{4}$. Και το $\frac{50}{60}=\frac{5}{6}$ (με το 10) και το $\frac{42}{60}=\frac{7}{10}$ (με το 6). Σωστά.

\tmima{Ασκήσεις}

\zoni{Το κατάλαβα}{εφαρμογή του κανόνα}

\ex Γράψε τρία κλάσματα ισοδύναμα με το $\frac{2}{3}$, πολλαπλασιάζοντας τους όρους του επί 2, επί 3 και επί 10.

\ex Βρες τον αριθμό που λείπει: α) $\frac{3}{4}=\frac{\square}{20}$ \quad β) $\frac{5}{7}=\frac{35}{\square}$ \quad γ) $\frac{18}{24}=\frac{3}{\square}$ \quad δ) $\frac{\square}{9}=\frac{20}{36}$.

\ex Έλεγξε με τα χιαστί γινόμενα ποια ζευγάρια είναι ισοδύναμα: α) $\frac{4}{6}$ και $\frac{10}{15}$ \quad β) $\frac{5}{8}$ και $\frac{15}{32}$ \quad γ) $\frac{7}{9}$ και $\frac{21}{27}$ \quad δ) $\frac{12}{18}$ και $\frac{14}{21}$ \quad ε) $\frac{3}{5}$ και $\frac{4}{7}$.

\ex Απλοποίησε μέχρι να γίνουν ανάγωγα: α) $\frac{18}{24}$ \quad β) $\frac{35}{42}$ \quad γ) $\frac{48}{72}$ \quad δ) $\frac{100}{250}$ \quad ε) $\frac{56}{98}$.

\ex Ποια από τα κλάσματα $\frac{9}{16}$, $\frac{12}{18}$, $\frac{15}{28}$, $\frac{21}{35}$, $\frac{17}{51}$ και $\frac{25}{36}$ είναι ανάγωγα; Για τα υπόλοιπα, γράψε το ανάγωγο.

\ex Μετάτρεψε σε ισοδύναμα κλάσματα με παρονομαστή 100: α) $\frac{3}{4}$ \quad β) $\frac{7}{20}$ \quad γ) $\frac{2}{5}$ \quad δ) $\frac{9}{50}$ \quad ε) $\frac{13}{25}$.

\ex Μετάτρεψε σε ισοδύναμα κλάσματα με παρονομαστή 12: α) $\frac{5}{6}$ \quad β) $\frac{3}{4}$ \quad γ) $\frac{2}{3}$ \quad δ) $\frac{1}{2}$.

\ex Κάνε ομώνυμα τα κλάσματα: α) $\frac{2}{3}$ και $\frac{3}{4}$ \quad β) $\frac{5}{6}$ και $\frac{7}{8}$ \quad γ) $\frac{4}{5}$, $\frac{1}{2}$ και $\frac{7}{10}$ \quad δ) $\frac{1}{6}$, $\frac{3}{8}$ και $\frac{5}{12}$.

\ex Σωστό ή λάθος; α) $\frac{2}{3}=\frac{4}{6}$ \quad β) $\frac{3}{5}=\frac{6}{12}$ \quad γ) Αν διαιρέσω τον αριθμητή με το 3 και τον παρονομαστή με το 4, βγαίνει ισοδύναμο κλάσμα. \quad δ) $\frac{7}{7}=1$ \quad ε) Το $\frac{6}{9}$ είναι ανάγωγο. \quad στ) Τα $\frac{4}{9}$ και $\frac{8}{18}$ είναι ισοδύναμα.

\zoni{Το χρησιμοποιώ}{προβλήματα από τη ζωή}

\ex \textbf{Οι μερίδες.} Τέσσερα παιδιά έφαγαν από μια ίδια πίτσα. Η Μυρτώ έφαγε $\frac{3}{4}$, η Φανή $\frac{6}{8}$, ο Ηλίας $\frac{9}{12}$ και ο Στέλιος $\frac{4}{6}$. Ποιος έφαγε περισσότερο και ποιος λιγότερο; Κάνε τα κλάσματα ομώνυμα για να βεβαιωθείς.

\ex \textbf{Η σοκολάτα.} Η σοκολάτα της Αγγελικής έχει 24 τετράγωνα και έφαγε τα 18. Τι κλάσμα της σοκολάτας έφαγε; Γράψε το ως ανάγωγο.

\ex \textbf{Τα λεπτά.} Η ώρα έχει 60 λεπτά. Γράψε ως ανάγωγο κλάσμα της ώρας τα 45, τα 20, τα 40 και τα 12 λεπτά.

\ex \textbf{Το τυρί της γιαγιάς.} Η γιαγιά Ρήνη αγοράζει τυρί 250, 750, 400 και 125 γραμμάρια σε διαφορετικές μέρες. Ποιο μέρος του κιλού (1.000 γραμμάρια) είναι κάθε ποσότητα; Γράψε το ανάγωγο κλάσμα.

\ex \textbf{Οι νότες της Αγγελικής.} Στο μάθημα της μουσικής η Αγγελική μαθαίνει ότι μια νότα «τέταρτο» διαρκεί το $\frac{1}{4}$ ενός μέτρου, μια «όγδοο» το $\frac{1}{8}$ και μια «δέκατο έκτο» το $\frac{1}{16}$. Πόσα όγδοα χωράνε σε ένα τέταρτο; Πόσα δέκατα έκτα; Γράψε τις ισότητες $\frac{1}{4}=\frac{\square}{8}=\frac{\square}{16}$.

\zoni{Το σκέφτομαι}{μια ερώτηση που πάει πιο μακριά}

\ex Ο Στέλιος έγραψε $\frac{1}{2}=\frac{1+3}{2+3}=\frac{4}{5}$. Δείξε με μια πίτσα ότι αυτό είναι λάθος. Ποιος είναι ο σωστός τρόπος να φτιάξεις ισοδύναμο κλάσμα με παρονομαστή 6;

\ex Πόσα διαφορετικά κλάσματα ισοδύναμα με το $\frac{3}{5}$ έχουν παρονομαστή μικρότερο από το 50; Γράψε τα.

\ex Βρες ένα κλάσμα ισοδύναμο με το $\frac{5}{8}$, του οποίου οι δύο όροι έχουν άθροισμα 91.

\ex Ένα κλάσμα είναι ισοδύναμο με το $\frac{3}{7}$ και οι όροι του διαφέρουν κατά 20. Ποιο είναι;

\ex Γιατί, όταν απλοποιούμε και τους δύο όρους με τον ΜΚΔ τους, το κλάσμα που βγαίνει είναι ανάγωγο; Σκέψου το $\frac{84}{126}$ και το $\frac{2}{3}$: αν τα 2 και 3 είχαν κοινό διαιρέτη μεγαλύτερο από το 1, τι θα συνέβαινε με το 42;

\ex Γράψε τρία κλάσματα που είναι ισοδύναμα με το 0 και τρία που είναι ισοδύναμα με το 1.

\zoni{Για το σπίτι}{με χαρτί και ψαλίδι}

\ex Δες τις τέσσερις ράβδους. Γράψε το κλάσμα του χρωματισμένου μέρους καθεμιάς και παρατήρησε τι συμβαίνει.\par\smallskip
\barfig{3}{2}\quad\barfig{6}{4}\quad\barfig{9}{6}\quad\barfig{12}{8}

\ex Πάρε ένα φύλλο χαρτί και δίπλωσέ το στη μέση, μετά πάλι στη μέση, και πάλι. Γράψε με κλάσματα το μέρος του φύλλου που έχει ένα στρώμα, δύο στρώματα και τέσσερα στρώματα όταν το ξεδιπλώσεις. Ποια ισοδύναμα κλάσματα βρήκες για το $\frac{1}{2}$;
